\bookchapter{Глава 62. Разработка интерфейсов — введение в SwiftUI}

\chapterintro{Интерфейс — это всё, что видит пользователь. SwiftUI — декларативный фреймворк интерфейсов, представленный Apple в 2019 году: вам нужно лишь «объявить», как выглядит интерфейс, а система сама займётся отрисовкой, тёмной темой и разными экранами. К 2026 году SwiftUI — безоговорочный мейнстрим для новых проектов. В этой главе мы научимся «рисовать».}

\booksection{Идея SwiftUI: декларативность, а не императивность}

Старая разработка интерфейсов — «императивная»: вы инструктируете систему шаг за шагом — «создай кнопку помести в (20,30) задай цвет». SwiftUI — «декларативный»: вы просто описываете «на экране есть красная кнопка с надписью “Нажми меня”», а как рисовать и когда перерисовывать — система решает сама.

\begin{lstlisting}[language=Swift]
import SwiftUI

struct ContentView: View {
    var body: some View {      // body: описывает «как выглядит экран»
        VStack {               // вертикально
            Text("Привет, мир!")
                .font(.largeTitle)      // модификатор: крупный заголовок
            Button("Нажми") {
                print("Нажато")
            }
            .buttonStyle(.borderedProminent)  // сплошная синяя кнопка
        }
        .padding()
    }
}
\end{lstlisting}

Ключевой момент: представление = структура, соблюдающая протокол View; свойство body описывает интерфейс, а модификаторы через точку «наносят макияж» цепочкой. Это и есть весь фундамент SwiftUI.

\booksection{Базовые представления: текст, кнопка, поле ввода, картинка}

\begin{center}
\begin{tabularx}{\linewidth}{|X|X|X|}
\hline
\textbf{Представление} & \textbf{Для чего} & \textbf{Пример} \tabularnewline
\hline
Text & показывать текст & Text("Привет") \tabularnewline
Button & кнопка & Button("Нажми меня") \{ действие \} \tabularnewline
TextField & поле ввода & TextField("подсказка", text: \$текст) \tabularnewline
Image & показывать картинку & \texttt{Image(systemName: "star.\allowbreak fill")} \tabularnewline
\hline
\end{tabularx}
\end{center}

\begin{center}\small\color{gray}Четыре базовых представления SwiftUI\end{center}

Обратите внимание на Image(systemName:) — у Apple есть встроенная библиотека иконок SF Symbols (несколько тысяч бесплатных иконок): имя вроде “star.fill” сразу покажет залитую звёздочку. Это скрытый бонус разработки под iOS: для интерфейсов почти не нужно искать иконки.

\booksection{Трио компоновки: VStack, HStack, ZStack}

\begin{itemize}
  \item VStack: выстраивает по вертикали (как очередь в столовой);
  \item HStack: выстраивает по горизонтали (как сидящие рядом);
  \item ZStack: накладывает друг на друга (как стопка игральных карт, подходит для бейджей и масок).
\end{itemize}

\begin{lstlisting}[language=Swift]
VStack(alignment: .leading, spacing: 12) {
    HStack {
        Text("Счёт")
        Spacer()                 // гибкое пространство: отталкивает правую часть к краю
        Text("-25.5 юаня").foregroundColor(.red)
    }
    Divider()                    // разделитель
    Text("Заметка: обед")
}
.padding(16)                     // отступы со всех сторон
.background(Color(.systemGray6)) // светло-серый фон
.cornerRadius(12)                // скругление углов

// Порядок модификаторов важен: сначала padding, потом background — фон не включает отступы
\end{lstlisting}

\begin{tcbwarn}{Внимание: порядок модификаторов имеет значение}
padding(16).background(серый) и background(серый).padding(16) дают разный эффект: первый — «сначала отступ, потом краска» (фон без отступов), второй — «сначала краска, потом отступ» (фон включает отступы). Отсюда и самые частые у новичков «странные белые полосы». Сомневаетесь — поменяйте порядок, и предпросмотр тут же покажет ответ.

\end{tcbwarn}

\booksection{Состояние и привязка: оживляем интерфейс}

Статичный интерфейс скучен. Правило SwiftUI: интерфейс всегда верно отражает состояние; изменилось состояние — интерфейс перерисовывается сам.

\begin{lstlisting}[language=Swift]
struct CounterView: View {
    @State private var count = 0   // @State: при изменении значения перерисовать

    var body: some View {
        VStack {
            Text("Вы нажали \(count) раз")
            Button("Нажми +1") {
                count += 1          // меняем состояние — интерфейс обновится сам
            }
        }
    }
}
\end{lstlisting}

Для поля ввода нужна «двусторонняя привязка» — знак \$ передаёт одновременно «значение и право его менять»:

\begin{lstlisting}[language=Swift]
struct InputView: View {
    @State private var text = ""

    var body: some View {
        TextField("Введите что-нибудь", text: $text)  // $ — двусторонняя связь
            .textFieldStyle(.roundedBorder)
        Text("Вы ввели: \(text)")   // обновляется на лету
    }
}
\end{lstlisting}

Действие пользователя

Состояние меняется (count += 1)

SwiftUI замечает изменение

Интерфейс перерисовывается сам (body пересчитывается)

\begin{center}\small\color{gray}Рисунок 3. Однонаправленный поток данных «состояние интерфейс» в SwiftUI\end{center}

\booksection{Предпросмотр, темы и анимация: «читы» SwiftUI}

\begin{itemize}
  \item Макрос \#Preview: добавьте в конец файла \#Preview \{ ContentView() \} — холст справа начнёт рендерить в реальном времени, и правки кода будут видны сразу;
  \item Тёмная тема бесплатно: используйте системные цвета (Color(.systemGray6), .primary), а не жёстко заданные — тёмная тема подстроится автоматически, без единой лишней строки;
  \item Анимация в одну строку: withAnimation \{ count += 1 \} или добавьте к представлению .animation(.default, value: count) — переходы станут шелковистыми.
\end{itemize}

\begin{lstlisting}[language=Swift]
struct AnimatedHeartView: View {
    @State private var liked = false

    var body: some View {
        Button {
            withAnimation(.spring(duration: 0.4)) {
                liked.toggle()
            }
        } label: {
            Image(systemName: liked ? "heart.fill" : "heart")
                .font(.system(size: 48))
                .foregroundColor(liked ? .red : .gray)
                .scaleEffect(liked ? 1.3 : 1.0)
        }
    }
}
\end{lstlisting}

\begin{tcbinfo}{Совет: сначала функция, потом анимация}
Типичная ошибка новичка — «интерфейс ещё не доделан, а он уже мучается с анимациями». Сначала добейтесь, чтобы всё работало, а эффекты добавьте в конце разом: анимация — украшение, функциональность — спасение. Это и есть наш вечный принцип «простота превыше всего».

\end{tcbinfo}

\begin{tcbpractice}{Попробуйте сами: сделайте «кнопку настроения»}
\booksection{Шпаргалка по частым модификаторам: одна страница под рукой — и интерфейсы не страшны}

\begin{center}
\begin{tabularx}{\linewidth}{|X|X|}
\hline
\textbf{Что хотите сделать} & \textbf{Модификатор} \tabularnewline
\hline
Отступы & \texttt{.\allowbreak padding(16)} \tabularnewline
Скругление углов & \texttt{.\allowbreak cornerRadius(12)} \tabularnewline
Цвет фона & \texttt{.\allowbreak background(Color(.\allowbreak systemGray6))} \tabularnewline
Шрифт и размер & \texttt{.\allowbreak font(.\allowbreak title.\allowbreak bold())}/\texttt{.\allowbreak font(.\allowbreak system(size: 20))} \tabularnewline
Цвет переднего плана (текст/иконки) & \texttt{.\allowbreak foregroundColor(.\allowbreak red)} \tabularnewline
Растянуть на всю ширину & \texttt{.\allowbreak frame(maxWidth: .\allowbreak infinity)} \tabularnewline
Фиксированный размер & \texttt{.\allowbreak frame(width: 100, height: 44)} \tabularnewline
Выравнивание & \texttt{VStack(alignment: .\allowbreak leading)} \tabularnewline
Прозрачность & \texttt{.\allowbreak opacity(0.\allowbreak 6)} \tabularnewline
Тень & \texttt{.\allowbreak shadow(radius: 5)} \tabularnewline
Жест нажатия & \texttt{.\allowbreak onTapGesture \{ .\allowbreak .\allowbreak .\allowbreak  \}} \tabularnewline
Тип клавиатуры & \texttt{.\allowbreak keyboardType(.\allowbreak decimalPad)} \tabularnewline
\hline
\end{tabularx}
\end{center}

\begin{center}\small\color{gray}Шпаргалка по частым модификаторам SwiftUI\end{center}

Ничего страшного, если не запомнили — эта страница и существует для справки. На деле чаще всего используются первые четыре: padding, скругление, фон, шрифт.

Переделайте шаблонный проект в переключатель настроений: одна кнопка, при каждом нажатии циклически меняется радость грусть злость, а текст и цвет фона меняются вместе с ней. Подсказка: храните три пары (текст + цвет) в массиве, индекс — в @State, при нажатии — (index + 1) \% 3. Когда сделаете, отправьте друзьям поиграть и соберите первую волну «вау, ты уже умеешь делать приложения!».

\end{tcbpractice}

\begin{tcbsummary}{Краткие итоги главы}
\begin{itemize}
  \item SwiftUI декларативен: вы описываете, «как это выглядит», а система рисует и перерисовывает
  \item представление = структура, соблюдающая протокол View; body описывает интерфейс; модификаторы «наносят макияж» цепочкой
  \item четыре базовых представления: Text / Button / TextField / Image; библиотека иконок SF Symbols — бесплатно и вдоволь
  \item трио компоновки: VStack — вертикаль, HStack — горизонталь, ZStack — наложение; Spacer — упругое пространство
  \item @State оживляет интерфейс; \$ — двусторонняя привязка для поля ввода
  \item \#Preview — что видишь, то и получаешь; системные цвета сами подстраиваются под тёмную тему; withAnimation — анимация в одну строку
\end{itemize}

\end{tcbsummary}
